\documentclass[11pt]{article}
\usepackage[T1]{fontenc}
\usepackage{lmodern}
\usepackage[a4paper,margin=25mm]{geometry}
\usepackage{amsmath,amssymb,amsthm,mathtools}
\usepackage{graphicx,booktabs,array,microtype}
\usepackage[font=small,labelfont=bf]{caption}
\usepackage[numbers,sort&compress]{natbib}
\usepackage{xcolor}
\usepackage[colorlinks=true,linkcolor=blue!50!black,citecolor=blue!50!black,%
            urlcolor=blue!50!black]{hyperref}

\newtheorem{theorem}{Theorem}[section]
\newtheorem{proposition}[theorem]{Proposition}
\newtheorem{lemma}[theorem]{Lemma}
\newtheorem{corollary}[theorem]{Corollary}
\theoremstyle{definition}
\newtheorem{remark}[theorem]{Remark}

\newcommand{\one}{\mathbf 1}
\newcommand{\diag}{\operatorname{diag}}
\newcommand{\TV}{\operatorname{TV}}
\newcommand{\E}{\mathbb E}
\newcommand{\Prob}{\mathbb P}
\newcommand{\norm}[1]{\left\lVert#1\right\rVert_\infty}
\newcommand{\R}{\mathbb R}
\newcommand{\kl}{\operatorname{kl}}
\setlength{\emergencystretch}{2em}

\title{Randomized stopped lineage experiments separate\\
phenotypic switching, selection and correlated inheritance}
\author{}
\date{20 September 2026}

\begin{document}
\maketitle

\begin{abstract}
The same resistant outgrowth can reflect pre-existing state, selective survival,
within-cell switching or altered daughter production. We ask which observations
separate these mechanisms in a finite-state continuous-time branching model with
correlated sister states. An independently randomized exponential deadline,
competing with the founder's first division or death, turns the observation law
into a resolvent: calibrated initial and endpoint markers together with joint
daughter records recover the switching generator, the demographic rates and the
daughter kernel at every finite rate and for any finite number of states,
without matrix logarithms and without a diagonalizability assumption. A
monitored delayed daughter readout preserves identification through an explicit
channel with interruption flags. For two states we give a rational formula for
the switching rate that cancels every marker determinant before it is evaluated,
and we replace Hoeffding radii by Chernoff--Kullback--Leibler feature sets, of
which the Hoeffding radii are exactly the Pinsker relaxation. On a synthetic
two-arm dataset with uncertain calibration these two changes narrow the
conservative $95\%$ induction interval from $[0.362,1.498]$ to $[0.578,1.092]$
per hour, around a true contrast of $0.8$ per hour, using the same observations
and the same coverage event; and they lower the certified budget of a
separated-class prospective decision from $2\times10^{6}$ to $3\times10^{5}$
founders. Known calibration mixtures replace pure-state calibration exactly, and
a six-category coarsening suffices when only within-cell switching is targeted.
Finally, increasing sister concordance at fixed daughter marginals increases
finite-time extinction weakly, and strictly in a stated case, while leaving
every mean trajectory unchanged. These are model-conditional observation and
prediction theorems, not a calibrated assay. The probabilistic arguments are
conventional proofs; a declared finite-algebra subset is verified in Lean.
\end{abstract}

\section{The question and its observation model}

Phenotypic inheritance, selection and induced plasticity are compatible
mechanisms. A lineage experiment should estimate their quantitative
contributions, rather than choose one of three mutually exclusive explanations.
Induction identification is already established in controlled ODE settings
\citep{greene2019}; branching inference distinguishes the information in counts
and fractions \citep{gunnarsson2023}, and division-associated plasticity is
explicitly modelled in drug-screen inference \citep{wu2025}. Lineage-tree
inference motivates tracking before treatment and after withdrawal, but a
lineage relationship is not itself a calibrated phenotype measurement
\citep{mohammadi2022}.

We give an observation-design result for a declared source. The distinguishing
step is to stop following a founder at its first demographic event, with an
independent random deadline. This exposes both a killed within-cell chain and
its joint daughter exit kernel. The resolvent method is classical in generator
estimation \citep{metzner2007}; stopped-process potential kernels are likewise
standard \citep{fitzsimmons1999}. Our claim is the explicit connection of those
tools to this lineage protocol, its delayed observation channel, its finite-data
target and its offspring-dependence consequence. We do not assert priority for
resolvent inversion or generating-function comparison in isolation.

Three questions are kept apart throughout. \emph{Structural identifiability}
asks whether a record law determines the source at all; it is answered by
algebra. \emph{Finite-sample inference} asks what a fixed number of independent
founders certifies; it is answered by concentration inequalities and exact
interval arithmetic. \emph{Biological calibration} asks whether a real reporter
measures the stipulated latent states; nothing in this paper answers it. Every
numerical value below is synthetic.

\subsection{Source, emissions and mechanisms}
The state space is $\{1,\ldots,m\}$. At a fixed exposure $u$, a cell of type $i$
switches to $j\ne i$ at rate $q_{ij}(u)$, dies at rate $d_i(u)$, and divides at
rate $b_i(u)$. Division replaces the mother by an ordered pair $(j,k)$ with
joint law $K_i^u(j,k)$. The kernel is stochastic; exchange symmetry is imposed
when the daughter order is arbitrary. Conditional on both realized daughter
states, their future subtrees have independent clocks. Daughter states need not
be independent conditional on the mother.

All per-cell rates are finite. The founder has probability column vector $\pi$,
with $\pi_i>0$. Experimental units are independent founder preparations in
controlled environments; treatment is randomized independently of founder state.
Known marker matrix $E\in\R^{p\times m}$ is column stochastic and has full
column rank. Its errors are independent conditional on the complete latent
history, and it does not perturb the source. Columns refer to the same
externally anchored states across arms. Distinct calibrated matrices at the
initial and subsequent reads are allowed; a common $E$ simplifies notation.
Rates and channels stay constant during founder and daughter follow-up.

Write $Q_{ii}=-\sum_{j\ne i}q_{ij}$, and define
\begin{equation}\label{eq:HB}
 H=Q-\diag(b+d),\qquad B=[\,d\mid (b_iK_i(j,k))_{i,(j,k)}\,].
\end{equation}
Thus $H\one+B\one=0$. Matrices act on row-state distributions. For two labelled
states $S,T$, the principal contrasts include
\begin{align}\label{eq:estimand}
 &\pi_T,\quad q_{ST}(1)-q_{ST}(0),\quad q_{TS}(0),\quad
 (b_T-d_T)-(b_S-d_S),\notag\\
 &L_T(T)-L_S(T),\quad c_i=K_i(T,T)-L_i(T)^2,\quad
 \zeta_i(u)=2b_i(u)L_i^u(T),
\end{align}
where $L_i(j)=\sum_kK_i(j,k)$ by exchange symmetry. The last contrast measures
the rate of tolerant daughter production at division; its treatment change is
distinct from within-cell induction. It is an expected output intensity, not a
probability of switching per division. The demographic contrast alone does not
describe composition dynamics when division changes phenotype.

Here ``predisposition'' means a pretreatment state predictive of a specified
later response. It does not mean an additional unobserved primed state. If
$b_i=0$, only $b_iK_i=0$ is identifiable: the unused kernel is immaterial. An
unanchored global label permutation is harmless; separate permutations in
different treatment arms are not.

\subsection{Existence and the branching law}
Independent competing exponential clocks, followed by kernel sampling at
division, construct the chronological source. If $b_i\le b_{\max}$, its
population is dominated by a Yule process of rate $b_{\max}$ per cell. The Yule
holding-time sum diverges almost surely: its centered independent terms have
summable variances, while their expectations form a divergent harmonic sum. On
every finite horizon the dominating population has finite integrated cell time,
so bounded switching and death intensities cannot accumulate infinitely many
additional events. The source is nonexplosive.

For $z\in[0,1]$ let $F_i(t,z)=\E_i z^{N(t)}$. First-event conditioning gives
\begin{equation}\label{eq:pgf}
 \partial_tF_i=d_i(1-F_i)+\sum_{j\ne i}q_{ij}(F_j-F_i)
       +b_i\left(\sum_{jk}K_i(j,k)F_jF_k-F_i\right),\quad F_i(0,z)=z.
\end{equation}
The product follows after conditioning on both daughter states. The bounded
first-event integral equation and uniqueness of the polynomial ODE in the unit
cube identify this with the actual source PGF. Extinction by time $t$ has
probability $F_i(t,0)$; extinct founders remain in the observed denominator.

\subsection{A precise limited-design obstruction}
\begin{proposition}\label{prop:alias}
Fix one treatment schedule and a time $t>0$, and retain only
$Y=\mathbf1\{N(t)>0\}$. Write $r_i=1-F_i(t,0)$ and $p=\pi_T$. At a two-state
source with $r_T\ne r_S$ and $0<p<1$, nearby changes in the rate/kernel
parameter $\theta$ can be compensated by
\begin{equation}\label{eq:alias}
 p'=\frac{(1-p)r_S+pr_T-r'_S}{r'_T-r'_S}.
\end{equation}
This is a local family on an observational equivalence level set. The full
binary observation law is unchanged, although within-cell induction, selection
or daughter inheritance may change.
\end{proposition}
\begin{proof}
The law is Bernoulli with parameter $(1-p)r_S+pr_T$. Continuous dependence in
\eqref{eq:pgf} keeps the denominator nonzero and $p'$ interior locally.
Substitution proves equality of the entire law. A nondegenerate class exists:
choose all rates positive and $d_S>d_T$. At zero, $(r_T-r_S)'=d_S-d_T>0$, so the
response contrast is positive for sufficiently small positive $t$. One may then
vary $q_{ST}(1)$ while holding the untreated rates fixed. Any test of two such
aliased sources has type-I plus type-II error one, since both product
observation laws are identical.
\end{proof}
The result concerns this binary threshold and an unknown preparation. It is not
an impossibility theorem for full fluctuation counts, all trees or multiple
randomized schedules.

\section{Randomized stopping and an algebraic source inverse}
\subsection{The retained record}
Read one founder marker $Y_0$ and draw an independent deadline
$T\sim\operatorname{Exp}(\lambda)$ with known $\lambda>0$. Follow the founder
until $\sigma=\min(T,\tau_{\rm dem})$, where $\tau_{\rm dem}$ is its first death
or division; unobserved switches do not stop follow-up. At a deadline, read an
endpoint marker. At death, retain a death flag. At division, read the joint
daughter marks. Neither the realized deadline nor the event timestamp enters the
retained record. All founders are retained. For binary markers this
immediate-read design has $2(2+1+4)=14$ outcomes. Figure~\ref{fig:assay}
summarises the protocol.

\begin{figure}[tb]
\centering\includegraphics[width=\linewidth]{figures/assay.pdf}
\caption{The observation protocol. The deadline clock is independent of the
source. Switching is hidden throughout. A delayed daughter read retains an
interruption flag for any intervening division or death; records are never
conditioned on both daughters surviving. The delayed binary-marker experiment
has 24 joint outcomes, versus 14 with immediate daughter readout.}
\label{fig:assay}
\end{figure}

\begin{theorem}[Finite-state randomized-deadline identification]\label{thm:resolvent}
Under the source and measurement assumptions above, put
$R_\lambda=(\lambda I-H)^{-1}$. The latent deadline and demographic-exit kernels
are
\begin{equation}\label{eq:reslaw}
 A_\lambda=\lambda R_\lambda,\qquad D_\lambda=R_\lambda B.
\end{equation}
The observed law determines $\pi,A_\lambda,D_\lambda$ and hence
\begin{equation}\label{eq:inverse}
 H=\lambda(I-A_\lambda^{-1}),\qquad B=\lambda A_\lambda^{-1}D_\lambda.
\end{equation}
It identifies $Q,b,d$ and each $K_i$ for which $b_i>0$, at arbitrary finite
rates and for any finite number of states. Reducibility, zero switching and
singular $H$ are allowed. No diagonalizability assumption is required.
\end{theorem}
\begin{proof}
Before demographic exit the founder is a killed chain: its row distribution
satisfies $v'=vH$, so its substochastic semigroup is $S_t=e^{tH}$. Given state
$i$ at time $t$, exit mark $a$ has intensity $B_{ia}$. Competition with the
independent deadline gives
\[
 A_\lambda=\int_0^\infty\lambda e^{-\lambda t}e^{tH}\,dt,
 \qquad D_\lambda=\int_0^\infty e^{-\lambda t}e^{tH}B\,dt.
\]
Since $\norm{e^{tH}}\le1$, these integrals converge. If $R$ denotes the integral
without $\lambda$ or $B$, integrate the derivative of $e^{t(H-\lambda I)}$ from
zero to infinity. Its limit is zero, and its derivative commutes with
$H-\lambda I$, giving both $(\lambda I-H)R=I$ and $R(\lambda I-H)=I$. This
proves the inverse and \eqref{eq:reslaw}, even if $H$ is singular. Moreover
\[
 A_\lambda\one+D_\lambda\one=R_\lambda(\lambda I-H)\one=\one,
\]
so death and division account for all probability missing from deadline records.

Let $C=\diag(1,E\otimes E)$ be the exit emission channel, with rectangular
blocks when necessary. For an arbitrary left inverse $LE=I$, a left inverse of
$C$ is $L_C=\diag(1,L\otimes L)$. The observed deadline and exit blocks are
\begin{equation}\label{eq:observed}
 J=E\diag(\pi)A_\lambda E^\top,\qquad V=E\diag(\pi)D_\lambda C^\top.
\end{equation}
These follow by summing over the initial, current and daughter birth states;
only measurement noise and daughter futures are factored, not the birth-state
kernel $K$. The marginal $m_0$ of the initial marker sums over every outcome, so
$\pi=Lm_0$. Thus
\[
 A_\lambda=\diag(\pi)^{-1}LJL^\top,\qquad
 D_\lambda=\diag(\pi)^{-1}LVL_C^\top.
\]
Positive founder masses permit row division. Equation~\eqref{eq:inverse} now
follows algebraically. Offdiagonal entries of $H$ give $q_{ij}$; the death column
of $B$ gives $d$, its daughter row sums give $b$, and division by positive $b_i$
gives $K_i$. The diagonal of $H$ checks the recovered total demographic killing.
Two randomized, label-aligned exposure arms identify both within-cell induction
and treatment changes in daughter production.
\end{proof}

\subsection{Finite horizons, stability and a model check}
The expected founder observation time is
\begin{equation}\label{eq:time}
 \E\sigma=\pi^\top R_\lambda\one\le\lambda^{-1}.
\end{equation}
For a physical cap $L_0$, use $\min(T,L_0)$ and label a cap endpoint as a
deadline endpoint. Coupling to the uncapped source shows
\begin{equation}\label{eq:cap}
 \TV(P_{\rm capped},P_{\rm uncapped})
 \le e^{-\lambda L_0}\pi^\top e^{L_0H}\one\le e^{-\lambda L_0}.
\end{equation}
Only paths with $T>L_0$ and no demographic event by $L_0$ can differ. A cap flag
could instead be retained for an exact capped likelihood. The capped kernel
itself is not a resolvent. Choosing $L_0=\log(1/\varepsilon)/\lambda$ yields a
feature-bias reserve $\varepsilon$.

If $\norm{H},\norm{H'}\le R_0$, then $\norm{A_\lambda^{-1}}\le1+R_0/\lambda$.
The inverse difference identity yields
\begin{equation}\label{eq:stability}
 \norm{H-H'}\le\lambda(1+R_0/\lambda)^2\norm{A_\lambda-A'_\lambda}.
\end{equation}
For binary markers, if every observed category differs by at most $e$ and
$\pi_i\ge\rho$, marker inversion and row normalization give
\[
 \norm{A_\lambda-A'_\lambda}
 \le \frac{2\norm{L}\norm{L^\top}+7\norm{L}}{\rho}\,e.
\]
Indeed the four-entry deadline block has maximum row-sum error $2e$; the initial
marginal has entrywise error at most $7e$; and $A'$ has row sums at most one.
Larger record spaces use their actual category count. The prefactor in
\eqref{eq:stability} is $\lambda+2R_0+R_0^2/\lambda$, minimized at
$\lambda=R_0$ when $R_0>0$. This is a conditioning guide, not an optimal design
theorem: event yield, feature variance and costs also depend on $\lambda$.

Two independently assigned deadline rates for the same source must obey
\begin{equation}\label{eq:identity}
 R_\lambda-R_\mu=(\mu-\lambda)R_\lambda R_\mu.
\end{equation}
This is the inverse difference identity, since both inverse matrices differ by
$(\mu-\lambda)I$. Recovered $H,B,\pi$ must agree across groups. Construct
simultaneous source confidence sets for each group and reject compatibility when
their intersection is empty. A union bound controls rejection of a true common
source. Passing finitely many checks does not establish a Markov model or
uniquely diagnose an incompatible biological assumption.

\subsection{Fixed deadlines as an alternative}
\begin{proposition}\label{prop:fixed}
For a deterministic deadline $h>0$, replace the kernels in \eqref{eq:reslaw} by
$A_h=e^{hH}$ and $D_h=W_hB$, $W_h=\int_0^h e^{tH}dt$. For two states the same
observations identify the source at every $h>0$. For general $m$, it suffices
that $h\norm{H}<\log2$ throughout the admitted class.
\end{proposition}
\begin{proof}
The killed-chain occupation derivation is unchanged. In two states write
$H=\left[\begin{smallmatrix}-\alpha&a\\r&-\beta\end{smallmatrix}\right]$. Its
eigenvalues are $[-(\alpha+\beta)\pm\sqrt{(\alpha-\beta)^2+4ar}]/2$, real and
nonpositive. If distinct, $A_h$ has distinct positive eigenvalues with the same
projectors; their unique real logarithms recover $H$. Any competing two-state
killed generator has real spectrum and commutes with $A_h$, so uses the same
projectors. In the repeated case $H=\xi I+N$, $N^2=0$, and $A_h=e^{h\xi}(I+hN)$
determines both $\xi$ and $N$, including $N=0$. The eigenvalues of $W_h$ are
$\int_0^h e^{t\xi}dt>0$; hence $B=W_h^{-1}D_h$. For general $m$ under the bound,
the convergent logarithm series on $\norm{A_h-I}<1$ gives $\log(e^{tH})=tH$ by
differentiating along $0\le t\le h$. Also $\norm{W_h/h-I}<1$, so a Neumann
series inverts $W_h$. Marker cancellation completes both cases. Randomized
stopping avoids these spectral and small-step restrictions.
\end{proof}

\subsection{Implementing a random deadline}\label{sec:implement}
Randomized stopping does not require physically interrupting a culture. For a
nonperturbing live reporter with sufficiently resolved complete traces, draw
$T_i\sim\operatorname{Exp}(\lambda)$ independently for each founder and retain
the record only up to $\min(T_i,\tau_{\mathrm{dem},i})$. The deadline may be
drawn before imaging, or afterwards provided it is drawn independently of the
recorded paths. It must not be chosen using the observed fate, and founders must
not be included only because their follow-up happened to be long enough.

Two finite-resource effects must be represented in the observation model rather
than assumed away. A finite movie horizon is the cap of \eqref{eq:cap}, or an
exact censored likelihood with a retained cap flag. A fixed camera frame
interval gives a discretized reading rule, not an exact continuously timed
marker read; the rounded-time channel must be derived or its discrepancy
bounded. Repeated randomized truncations of one founder trace remain dependent
observations and cannot be counted as additional independent roots.

Raw times and images should be retained when feasible, even though
Theorem~\ref{thm:resolvent} shows they need not enter the minimal statistical
record. Compression proves sufficiency for structural identification on this
model; it does not imply that timestamps carry no additional finite-sample
information. Randomization creates no information that was absent from the
original trace.

\section{Delayed daughter observations with interruption flags}
After a founder division, follow each daughter for a known duration $\tau$. If
she divides or dies during that interval, retain an interruption flag and stop
following her. Otherwise read her marker at $\tau$. Switching during the
interval remains hidden. The two daughter records are retained jointly,
including all interrupted cases. No inference conditions on joint survival.

\begin{theorem}\label{thm:delay}
For every finite $\tau\ge0$, the daughter observation channel is
\begin{equation}\label{eq:delay}
 F_\tau=\begin{bmatrix}E e^{\tau H^\top}\\
                  (\one-e^{\tau H}\one)^\top\end{bmatrix}.
\end{equation}
Its columns sum to one and it has full column rank. The monitored delayed
protocol identifies the same source as Theorem~\ref{thm:resolvent}, or
Proposition~\ref{prop:fixed} within its scope.
\end{theorem}
\begin{proof}
The unnormalized live-state distribution of a daughter born in state $j$ is row
$j$ of $e^{\tau H}$. Applying $E$ gives the top block of \eqref{eq:delay}; its
missing mass is the interruption probability. Since $e^{\tau H}$ is invertible
and $E$ has full column rank, the upper block has full column rank. An explicit
left inverse is $e^{-\tau H^\top}L[\,I_p\mid0\,]$.

The founder deadline block, unaffected by daughter readout, first recovers
$\pi,H$. Summing delayed pair categories recovers founder division totals; the
two demographic columns are $R_\lambda[d\mid b]$, so recover $d,b$. Now $H$
determines $F_\tau$. Conditional on the two daughter birth states, the two
delayed records are independent; hence their joint channel is
$F_\tau\otimes F_\tau$. If $L_\tau F_\tau=I$, its left inverse is
$L_\tau\otimes L_\tau$, as follows by multiplying and factoring the two finite
sums. Substitute this channel in \eqref{eq:observed} to recover $b_iK_i(j,k)$,
and then $K_i$ wherever $b_i>0$. This ordering avoids circular inference of $H$
from an unknown daughter channel.
\end{proof}
The binary-marker record has $2(2+1+9)=24$ categories. A full-law Hoeffding
region over two arms therefore has radius $\sqrt{\log(96/\delta)/(2n)}$, rather
than the immediate-record constant. The direct switching interval of
Section~\ref{sec:inference} uses only its ten stated features and does not
require all 48 categories to be estimated separately.

Invertibility is not good conditioning: in spectral norm,
$\sigma_{\min}(F_\tau\otimes F_\tau)=\sigma_{\min}(F_\tau)^2$. If the founder
division probability is $y_\lambda=\pi^\top R_\lambda b$, additional daughter
observation is at most $2\tau y_\lambda$ cell-hours per founder. Equal root
counts and equal observation-time budgets are different comparisons
(Figure~\ref{fig:delay}). Monitoring still must detect the first founder
division and daughter interruptions. Unknown reporter maturation, inherited
fluorescence, unknown delay or missed events are not covered by
\eqref{eq:delay}.

Discarding the second daughter loses genuine information: the perturbation
$K_i\mapsto K_i+\varepsilon(1,-1,-1,1)$ preserves each daughter marginal but
changes $c_i$ by $\varepsilon$. For interior kernels small nonzero $\varepsilon$
is feasible. The whole one-daughter stopped law remains equal, whereas the
invertible pair channel separates the joint kernels. Thus the requirement of
joint daughter records is substantive.

\section{Induction intervals with unknown nuisance rates}\label{sec:inference}
The exact inverse need not be applied to every parameter in order to estimate a
particular mechanism. For two states let $M=\diag(\pi)A_\lambda=LJL^\top$.
Direct inversion of a two-by-two matrix gives
\begin{equation}\label{eq:target}
 q_{ST}=\lambda\frac{(A_\lambda)_{ST}}{\det A_\lambda}
       =\lambda\frac{\pi_T M_{ST}}{\det M},\qquad
 q_{TS}=\lambda\frac{\pi_S M_{TS}}{\det M}.
\end{equation}
The factors $\pi_T$ and $\pi_S$ follow from the row-state convention. The
determinant is positive: $\lambda I-H$ has positive diagonal entries and
determinant $(\lambda+a+b_S+d_S)(\lambda+r+b_T+d_T)-ar>0$. Consequently
$\det A_\lambda>0$ and $\det M>0$ for positive founder masses.

Only five features per arm are needed: the four joint initial/deadline marker
probabilities and $\Prob(Y_0=1)$ over \emph{all} founders. Demographic rates and
$K$ influence these probabilities but need not be estimated or assumed known.
The same formula works with delayed daughter records, since none of these five
features uses daughter outcomes. We therefore say the construction proceeds
\emph{without estimating demographic nuisance rates}; calibration, founder
prevalence and the source model still matter, and $q_{ST}$ is a switching rate
whereas induction is the between-arm contrast.

\subsection{Cancelling the marker determinant}\label{sec:cancel}
Evaluating \eqref{eq:target} by first widening each entry of $M$ and then taking
a difference of products repeats the same uncertain calibration factors several
times. They can be cancelled exactly beforehand.

\begin{proposition}[Cancelled-determinant target]\label{prop:cancel}
Write $e_S=\Prob(Y=1\mid S)$, $e_T=\Prob(Y=1\mid T)$, $\delta_E=e_T-e_S>0$,
$\mu=\Prob(Y_0=1)$ and $J=(j_{ab})_{a,b\in\{0,1\}}$. Put
\begin{equation}\label{eq:N}
 N=-e_Te_Sj_{00}+e_T(1-e_S)j_{01}+(1-e_T)e_Sj_{10}-(1-e_T)(1-e_S)j_{11}.
\end{equation}
Then $M_{ST}=N/\delta_E^2$, $\det M=\det J/\delta_E^{2}$ and
$\pi_T=(\mu-e_S)/\delta_E$, so that
\begin{equation}\label{eq:S1}
 \boxed{\;q_{ST}
  =\lambda\,\frac{(\mu-e_S)\,N}{\delta_E\,(j_{00}j_{11}-j_{01}j_{10})}.\;}
\end{equation}
\end{proposition}
\begin{proof}
With
\begin{equation}\label{eq:L}
 L=\frac1{\delta_E}\begin{bmatrix}e_T&-(1-e_T)\\-e_S&1-e_S\end{bmatrix},
 \qquad \pi_T=\frac{\mu-e_S}{\delta_E},
\end{equation}
one has $\det L=\bigl(e_T(1-e_S)-(1-e_T)e_S\bigr)/\delta_E^2=1/\delta_E$. Hence
$\det M=\det(LJL^\top)=(\det L)^2\det J=\det J/\delta_E^2$. Expanding
$M_{ST}=(LJL^\top)_{01}$ gives exactly $N/\delta_E^2$. Substituting both into
\eqref{eq:target} cancels $\delta_E^{-2}$ and leaves \eqref{eq:S1}.
\end{proof}

Identity \eqref{eq:S1} is exact, so it changes no modelling assumption. Its
value is that in interval arithmetic every factor appears as few times as
possible: the calibration contrast $\delta_E$ occurs once instead of four times,
and the four observed probabilities enter one bilinear numerator and one
determinant instead of four separately widened entries of $M$. We do not claim a
universal width inequality between the two evaluations; we report the measured
reduction on the dataset of Section~\ref{sec:worked} and across the matched
diagnostics of Figure~\ref{fig:inference}.

\subsection{Chernoff feature sets}\label{sec:chernoff}
Let $\kl(a\Vert p)=a\log\frac ap+(1-a)\log\frac{1-a}{1-p}$ be the binary
relative entropy, with the usual conventions at $a\in\{0,1\}$.

\begin{lemma}[Chernoff set and its Pinsker relaxation]\label{lem:chernoff}
Let $\hat p$ be the mean of $n$ independent Bernoulli$(p)$ variables and $L>0$.
Put $C_L(\hat p)=\{q\in(0,1):n\,\kl(\hat p\Vert q)\le L\}$. Then
$\Prob\bigl(p\notin C_L(\hat p)\bigr)\le2e^{-L}$, and
\begin{equation}\label{eq:pinsker}
  C_L(\hat p)\subseteq\Bigl[\hat p-\sqrt{L/(2n)},\;\hat p+\sqrt{L/(2n)}\Bigr].
\end{equation}
\end{lemma}
\begin{proof}
The map $a\mapsto\kl(a\Vert p)$ is convex with minimum $0$ at $a=p$, so
$\{p\notin C_L(\hat p)\}=\{\hat p\ge a_+\}\cup\{\hat p\le a_-\}$, where
$\kl(a_\pm\Vert p)=L/n$ and $a_-\le p\le a_+$ (an empty event if the
corresponding root does not exist in $[0,1]$). The Chernoff bound for binomial
tails \citep{hoeffding1963,boucheron2013} gives
$\Prob(\hat p\ge a)\le e^{-n\kl(a\Vert p)}$ for $a\ge p$ and
$\Prob(\hat p\le a)\le e^{-n\kl(a\Vert p)}$ for $a\le p$; each term is at most
$e^{-L}$. Pinsker's inequality in the binary case is
$\kl(a\Vert q)\ge2(a-q)^2$, so $q\in C_L(\hat p)$ forces
$2(\hat p-q)^2\le L/n$, which is \eqref{eq:pinsker}.
\end{proof}

Taking $L=\tfrac{25}4$ recovers exactly the radius $r_n=\sqrt{25/(8n)}$, and
$L=6$ recovers exactly $s_c=\sqrt{3/n_c}$. In other words, the explicit
Hoeffding radii are the Pinsker relaxation of the Chernoff sets: the sets have
the same error allocation and are never wider, strictly narrower whenever the
proportion is away from one half. Such $\kl$-based intervals are standard in the
concentration literature \citep{garivier2011}. Because rational interval
arithmetic is inclusion preserving, we obtain the following free improvement.

\begin{corollary}\label{cor:never-worse}
Fix any evaluation of \eqref{eq:target} or \eqref{eq:S1} by inclusion-preserving
interval arithmetic. Replacing each Hoeffding feature interval by the
corresponding Chernoff set $C_L$ yields an output interval contained in the
original output interval, with an unchanged coverage guarantee.
\end{corollary}

\subsection{The finite-sample interval}
\begin{theorem}[Direct finite-sample interval]\label{thm:CI}
Take $n$ independent founders per arm and $n_c$ independent labelled calibration
readouts per state, independent of the experimental units. Replace each of the
ten empirical feature proportions by the Chernoff set $C_{25/4}$ of
Lemma~\ref{lem:chernoff}, and each of the two calibration proportions by
$C_{6}$. Evaluate \eqref{eq:L} and \eqref{eq:S1} using inclusion-preserving
interval arithmetic. If the marker-contrast lower bound or the lower bound on
$\det J$ is nonpositive, return $[0,\infty)$ for that arm; otherwise intersect
the result with $[0,\infty)$, an empty result being a model-conflict flag.
Subtract the arm intervals to obtain an induction interval. Its coverage is at
least $0.95$, uniformly over the admitted source, without knowing division,
death or reverse-switching rates, $K$ or $\pi$.
\end{theorem}
\begin{proof}
By Lemma~\ref{lem:chernoff} each founder feature is missed with probability at
most $2e^{-25/4}$ and each calibration proportion with probability at most
$2e^{-6}$. The ten features may overlap within a unit; the union bound does not
require their independence. Hence all twelve true values lie in their sets
except with probability at most $20e^{-25/4}+4e^{-6}<0.04+0.01=0.05$. On that
joint coverage event, rational interval addition, multiplication and division by
a sign-separated interval contain the true intermediate values, by their
endpoint definitions. The true $\delta_E$ and $\det J$ are positive, the
latter because $\det J=\delta_E^2\det M$ and $\det M>0$; when the enclosures
certify those two signs the evaluation proceeds, and otherwise the unbounded
fallback is returned. Thus
\eqref{eq:S1}, which by Proposition~\ref{prop:cancel} equals $q_{ST}$ exactly,
is contained whenever it is evaluated, and the fallback contains it
otherwise. The nonnegative intersection cannot remove the true switching rate.
Subtracting the intervals retains the true contrast. No asymptotic likelihood
argument and no uncertified global optimization is used.
\end{proof}

The implementation encloses each Chernoff set by certified rational endpoints:
candidate endpoints are located in floating point and then \emph{verified} to lie
outside the set by an exact rational lower bound on $\kl$, built from an
argument-reduced arctangent-hyperbolic series with an explicit geometric tail
(Appendix~\ref{app:numerics}). The inequalities $e^6>400$ and $e^{25/4}>500$
follow from finite positive Taylor sums. Arithmetic guarantees containment;
concentration together with the independence and model assumptions supplies the
probability guarantee. Both are needed.

The exact binomial (Clopper--Pearson) set is narrower still, and it is the
natural next candidate. We do not use it, for an implementation reason that is
specific to an exact-arithmetic pipeline: certifying a Clopper--Pearson endpoint
$p$ requires bounding a partial binomial tail $\sum_{k\le x}\binom nk
p^k(1-p)^{n-k}$, and at $n=5\times10^4$ a rational $p$ with a $10^{-12}$ grid
makes $p^k$ carry a denominator of order $10^{12k}$, so the exact sum is not
computable. The Chernoff set is the tightest object we know how to certify in
exact rational arithmetic, because its defining inequality involves only two
logarithms, each of which has a short rational enclosure. A floating-point
Clopper--Pearson interval would reintroduce exactly the uncertified numerics
that the rest of this construction avoids.

For a cap or other justified per-feature bias $\eta$, widen each feature
interval by $\eta$. Channel calibration uncertainty is already carried as a
region for $(e_S,e_T)$ and is not silently replaced by a point estimate. With
known $E$, omit the calibration error allocation. If reporter errors depend on
unobserved histories rather than the stipulated current state, a separate
observation model or quantified error bound is necessary.

\subsection{A worked unknown-nuisance dataset}\label{sec:worked}
All quantities in this example are synthetic. The generating source uses
\begin{gather}
 \pi_T=.35,\quad b=(.22,.16),\quad d=(.12,.04),\quad q_{TS}=.07,\notag\\
 K_S=(.62,.12,.12,.14),\quad K_T=(.10,.10,.10,.70),\notag\\
 E=\begin{bmatrix}.9&.15\\.1&.85\end{bmatrix},\quad
 q_{ST}(0)=.05,\quad q_{ST}(1)=.85,\quad\lambda=1.\label{eq:example}
\end{gather}
Rates are per hour. The inference routine receives only category counts,
calibration counts and $\lambda$, not these source parameters. For $n=50{,}000$
per arm and $n_c=100{,}000$ per state, seed 20260922 gives calibration successes
$(10{,}211,85{,}049)$ and the counts of Table~\ref{tab:counts}.

\begin{table}[tb]
\caption{The retained two-arm dataset. $J$ columns refer to the endpoint marker
only when the deadline precedes demographic exit; the fifth column also includes
death and division records. The last two columns are the demographic-exit totals
of the six-category coarsening of Section~\ref{sec:six}.}
\label{tab:counts}\centering\small
\begin{tabular}{lrrrrrrr}
\toprule
& \multicolumn{4}{c}{deadline first} & & \multicolumn{2}{c}{demographic exit}\\
\cmidrule(lr){2-5}\cmidrule(lr){7-8}
Arm & $J_{00}$ & $J_{01}$ & $J_{10}$ & $J_{11}$ & all initial $Y_0=1$
    & $Y_0=0$ & $Y_0=1$\\
\midrule
Untreated & 19497 & 4654 & 4270 & 10288 & 18023 & 7826 & 3465\\
Treated   & 13937 & 10895 & 3754 & 11178 & 18242 & 6926 & 3310\\
\bottomrule
\end{tabular}
\end{table}

Table~\ref{tab:routes} evaluates the same data and the same coverage event by
four routes. The original construction --- Hoeffding radii and entrywise
inversion of $M$ --- returns the induction interval $[0.362,1.498]$. Cancelling
the marker determinant removes $39.47\%$ of the width; replacing the radii by
Chernoff sets removes $30.58\%$; together they remove $54.79\%$. The
outward-rounded publication interval is
\begin{equation}\label{eq:workedCI}
                 [0.578,\,1.092]\ \mathrm{hour}^{-1},
\end{equation}
around a true contrast of $0.8$ per hour. It excludes zero. The improvement uses
the same observations, the same error allocation and the same confidence event
as the original construction; it introduces no new statistical assumption, extra
data, asymptotic approximation or additional error allocation, and it does not
by itself establish a universal sample-size requirement.

\begin{table}[tb]
\caption{Four evaluations of the same data under the same $95\%$ coverage event.
Endpoints are outward-rounded to four decimals; the exact rational
endpoints are in the accompanying numerical record. The
Chernoff rows are contained in the Hoeffding rows of the same evaluator
(Corollary~\ref{cor:never-worse}); the cancelled columns are narrower here but
are not claimed to be narrower at every dataset.}
\label{tab:routes}\centering\small
\begin{tabular}{llccc r}
\toprule
Feature sets & Evaluator & Untreated $q_{ST}$ & Treated $q_{ST}$ & Induction &
Width reduction\\
\midrule
Hoeffding & \eqref{eq:target} & $[0,.1608]$ & $[.5232,1.4978]$ & $[.3624,1.4978]$ & ---\\
Hoeffding & \eqref{eq:S1}     & $[0,.1282]$ & $[.6316,1.1908]$ & $[.5034,1.1908]$ & $39.47\%$\\
Chernoff  & \eqref{eq:target} & $[.0054,.1190]$ & $[.5999,1.2746]$ & $[.4809,1.2692]$ & $30.58\%$\\
Chernoff  & \eqref{eq:S1}     & $[.0130,.1010]$ & $[.6793,1.1045]$ & $[.5783,1.0916]$ & $54.79\%$\\
\bottomrule
\end{tabular}
\end{table}

A qualitative gain accompanies the quantitative one. Under the original route
the untreated arm returned $[0,.161]$, consistent with no spontaneous switching
at all; the improved route returns $[.0130,.1010]$, which certifies a strictly
positive baseline switching rate at the same confidence level. The true value is
$.05$ per hour.

A retained 20-founder example has a determinant interval reaching zero and
returns unresolved under all four routes. The procedure never replaces that
failure with a clipped finite answer.

In 300 independent synthetic repetitions of the improved route, including fresh
calibration each time, the interval covered the true contrast and excluded zero
in all 300. For either empirical proportion the two-sided $95\%$ exact binomial
Monte Carlo interval is $[.9878,1]$. This is a diagnostic of conservatism and
power at this source, not the coverage proof and not a universal power claim.

\begin{figure}[tb]
\centering\includegraphics[width=\linewidth]{figures/inference.pdf}
\caption{Matched design diagnostics for the sources in \eqref{eq:example}. All
panels use the same source pair and 200,000 total calibration readouts, rounded
expected category counts and the conservative interval algorithm; they are not
Monte Carlo coverage or prospective guarantees, and unbounded intervals are
shown as unresolved rather than as finite widths. Left: varying the number of
founders in each arm at four deadline rates, using the improved route. Middle:
varying total expected founder observation time with equal allocation to the two
arms; calibration time is not included in founder-hours and its sample count is
held fixed. Right: the four evaluation routes of Table~\ref{tab:routes} at
$\lambda=1$.}
\label{fig:inference}
\end{figure}

\subsection{A switching-only record with six categories}\label{sec:six}
If the scientific target is only within-cell switching, the daughter phenotype
readout can be dropped entirely.

\begin{proposition}\label{prop:six}
Retain, per arm, the four deadline categories $(Y_0,Y_T)$ and two
demographic-exit categories indexed by $Y_0$ alone, with counts
$(n_{00},n_{01},n_{10},n_{11},n_{0D},n_{1D})$. These six categories determine
all five features of \eqref{eq:S1}, hence Theorem~\ref{thm:CI} applies verbatim
to this record.
\end{proposition}
\begin{proof}
The four deadline proportions are immediate. The fifth feature is
$\mu=(n_{10}+n_{11}+n_{1D})/n$, and $n_{0D}$ is determined by subtraction from
$n$. No other functional of the record enters \eqref{eq:S1}.
\end{proof}

Here $D$ means \emph{either} first death \emph{or} first division, for this
deliberately coarsened target; it does not mean death alone. The corresponding
counts for the worked dataset are the last two columns of
Table~\ref{tab:counts}. This six-category multinomial retains all five features
including their dependencies, so no information about $q_{ST}$ is lost relative
to the full record. Recovering $b,d,K$, or predicting the sister-dependence
consequence of Section~\ref{sec:extinction}, still requires the richer record.
This is a useful experimental cost distinction, not a claim that all mechanisms
can be recovered from six categories.

\subsection{Prospective guarantees and the source of cost}
A prospective statement must control every source in a specified class, not just
the realized dataset. Here is a deliberately local, nonempty example. Allow
$\pi_T\in[.345,.355]$, $q_{TS}\in[.068,.072]$, $b_S\in[.218,.222]$,
$b_T\in[.158,.162]$, $d_S\in[.118,.122]$, $d_T\in[.038,.042]$, and arbitrary
stochastic $K$ in each arm. Set the true marker probabilities to $.1,.85$, but
estimate them by the calibration procedure. Control $q_{ST}$ lies in
$[.049,.051]$; the null treated range is the same, whereas the alternative
treated range is $[.849,.851]$. Thus neither hypothesis hides a known
demographic nuisance value, although these performance classes are intentionally
narrow. Classify using the threshold $.4$ per hour.

For this calculation, enclose the exact two-state inverse $(\lambda I-H)^{-1}$
over the input boxes and apply the emissions, obtaining a box of true feature
values. Then enlarge that box to cover every interval the procedure could
report. Under the original construction the enlargement is by $2r_n$ and $2s_c$,
the factor two accounting for sampling displacement and the reported confidence
radius. Under Chernoff sets the corresponding enlargement of a true value $p$ is
\begin{equation}\label{eq:double}
 \mathcal E(p)=\bigcup\{C_L(a):\;n\,\kl(a\Vert p)\le L\}
   =\bigl[\,\ell\bigl(a_-(p)\bigr),\;u\bigl(a_+(p)\bigr)\,\bigr],
\end{equation}
where $\ell,u$ are the endpoints of $C_L$ and $a_\pm$ those of the ball in the
first argument; both maps are nondecreasing, so the union is the stated
interval and a box of true values is enlarged by applying $\ell$ and $u$ at its
two ends. By \eqref{eq:pinsker}, $\mathcal E(p)\subseteq[p-2r_n,p+2r_n]$, so the
prospective enclosure is never wider than before. Inclusion arithmetic encloses
every possible output on the coverage event; no source grid or optimizer
establishes the extrema.

At $\lambda=1$ the original route certifies the decision at
$n=n_c=10^{6}$ per arm and per calibration state: all returned null contrast
intervals lie inside $[-.129,.129]$ and all returned alternative intervals
inside $[.441,1.299]$. The improved route certifies the same decision at
$n=n_c=1.5\times10^{5}$, with null intervals inside $[-.144,.144]$ and
alternative intervals inside $[.408,1.314]$; it does not certify at
$n=n_c=1.4\times10^{5}$, where the alternative lower bound is $.3998$. The
sufficient budget therefore falls from two million founders plus two million
labelled calibration readouts to three hundred thousand of each, a factor of
$6.7$. These remain sufficient budgets for this narrow class, not recommended
experimental sizes and not lower bounds. At smaller budgets the rectangular
method still loses information through interval dependencies, marker subtraction
and determinant conditioning.

There is nevertheless an intrinsic small-effect limit for the immediate-read
record, or its founder-only coarsening. Compare sources differing only in
$q_{ST}=a$ versus $a+\epsilon$, $a\ge a_0>0$. The latent stopped-path relative
entropy is expected occupation in $S$ times
$a\log(a/(a+\epsilon))+\epsilon\le\epsilon^2/a_0$. Expected occupation is at
most $1/\lambda$; discarding latent information cannot increase relative
entropy. Consequently any test on $n$ treated founders has error sum at least
\begin{equation}\label{eq:lower}
 \max\left\{0,1-\sqrt{\frac{n\epsilon^2}{2a_0\lambda}}\right\}.
\end{equation}
This follows from the variational testing bound, Pinsker's inequality and
product additivity; the unchanged control adds no information about this pair.
It explains why arbitrarily small induction cannot be resolved uniformly at a
fixed sample size. It does not establish that the sufficient budgets above are
necessary.

\subsection{What a joint region could still buy}\label{sec:joint}
Rectangular propagation discards the dependence between the features, and a
better region is plausible. It is worth stating precisely what would be needed,
because a Dirichlet posterior credible set is not automatically a uniformly
valid frequentist confidence set, and an ordinary likelihood ellipsoid is
generally only asymptotic.

\begin{proposition}[Normalized-mixture region]\label{prop:mixture}
For one ordered categorical sample with counts $n_1,\ldots,n_k$ let
$\mathcal L(p;x)=\prod_rp_r^{n_r}$ and $m(x)=\int\mathcal L(v;x)\,\Pi(dv)$, where
the mixing distribution $\Pi$ on the simplex is fixed before the data are seen.
Then $\mathcal C_\alpha(x)=\{p:\mathcal L(p;x)\ge\alpha\,m(x)\}$ has coverage at
least $1-\alpha$ for every $p$. For the uniform Dirichlet mixture,
$m(x)=(k-1)!\prod_rn_r!/(n+k-1)!$.
\end{proposition}
\begin{proof}
For fixed $p$, $\mathcal L(p;\cdot)$ is a probability distribution on ordered
samples of length $n$, so $\sum_x m(x)=\int\sum_x\mathcal L(v;x)\Pi(dv)=1$ by
Fubini. Hence $\E_p[m(X)/\mathcal L(p;X)]=\sum_{x:\mathcal L(p;x)>0}m(x)\le1$,
where samples of zero probability contribute nothing. Markov's inequality gives
$\Prob_p(m(X)/\mathcal L(p;X)>1/\alpha)\le\alpha$, which is the complement of
$\{p\in\mathcal C_\alpha(X)\}$. The Dirichlet integral is the Beta function
identity.
\end{proof}

Multiplying the mixture distributions for the two independent founder samples
and the two independent calibration samples gives one normalized joint
numerator, so a single threshold $\alpha=.05$ delivers simultaneous coverage
without splitting the error budget across twelve marginal sets. Intersecting the
region with the calibrated common-label and valid-source constraints and then
projecting the induction target preserves coverage, provided the projection is
outer-validated; an inner feasible region alone does not, because it can exclude
the truth. The exact arithmetic numerator is straightforward, and certified
projection of the nonlinear target is the real computational task; it is not
implemented here. A closely related established alternative is the split
likelihood-ratio construction of universal inference \citep{wasserman2020},
which supplies finite-sample likelihood-ratio confidence sets and nuisance
handling; its precise construction should be used rather than attaching a
finite-sample label to an ordinary likelihood ellipsoid.

Mixture penalties can themselves be conservative, so superiority over the
simplified rectangle is not promised. The improvements actually carried out in
this paper --- Proposition~\ref{prop:cancel} and Lemma~\ref{lem:chernoff} ---
are both exact and both reuse the original coverage event, which is why they
compose without any further union bound. We also note that one must not select
the most favourable among several separately valid procedures without accounting
for the selection; intersecting enclosures derived from the \emph{same}
simultaneous event, as here, is a different and legitimate operation.

\section{Calibration without pure populations}\label{sec:calibration}
Theorem~\ref{thm:CI} requires independently justified state labels in the
calibration sample. Pure sorted populations are a strong demand, and more
calibration replicates reduce sampling uncertainty but never unknown systematic
label contamination. Known mixtures suffice instead, exactly.

\begin{proposition}[Known-mixture calibration]\label{prop:mixcal}
Let two calibration preparations have independently known tolerant fractions
$p_0\ne p_1$ \emph{at readout}, and let $r_g=\Prob(Y=1\mid\text{preparation }g)$.
Under the same state-conditional channel, $r_g=e_S+(e_T-e_S)p_g$, hence
\begin{equation}\label{eq:S2}
 e_T-e_S=\frac{r_1-r_0}{p_1-p_0},\quad
 e_S=\frac{p_1r_0-p_0r_1}{p_1-p_0},\quad
 e_T=\frac{(1-p_0)r_1-(1-p_1)r_0}{p_1-p_0}.
\end{equation}
Pure states are the special case $p_0=0$, $p_1=1$. More generally, if $P$ is a
known matrix of preparation compositions with full row rank and $Y=EP$ are the
observed marker distributions, any right inverse gives $E=YP^{+}$.
\end{proposition}
\begin{proof}
Condition on the latent state within a preparation and use that the channel does
not depend on the preparation. The displayed formulas solve the resulting linear
system, whose determinant is $p_1-p_0$. The matrix statement is the same
computation with $P$ in place of $(p_0,p_1)$.
\end{proof}

With uncertain mixture fractions, propagate their independently justified
confidence or bias regions jointly with $r_0,r_1$ and enforce a nonzero mixture
contrast; the additional factor $|p_1-p_0|^{-1}$ quantifies the price of similar
preparations. This does not solve unknown ground truth by itself: the
compositions must be established by information independent of the uncalibrated
marker being assessed. Sorting on that same marker and declaring the gates pure
is circular.

\begin{remark}[Bounded contamination]\label{rem:contam}
If nominal $S$ and $T$ preparations have contamination fractions at most
$\eta_S,\eta_T$ at readout, their marker means differ from $e_S,e_T$ by at most
$\eta_S|e_T-e_S|$ and $\eta_T|e_T-e_S|$, hence by at most $\eta_S$ and $\eta_T$.
Widen the calibration sets by these justified bounds and propagate the resulting
channel region. Unknown contamination must not be converted into an arbitrarily
small reserve. If cells switch between sorting and measurement, the relevant
composition is the composition at readout; a known transition model and delay
can propagate it, and without that knowledge sorting purity at time zero does not
establish readout purity. If the resulting channel region admits rank loss,
report partial or unresolved identification.
\end{remark}

\section{Sister concordance changes extinction without changing means}
\label{sec:extinction}
The joint kernel is identifiable for a reason beyond fitting lineage
correlations: it determines a population outcome that mean growth cannot.
Offspring generating-function orders have an established role in extinction
theory \citep{bertacchi2026}. The following finite-time consequence is
specialized to the identified two-state source and a fixed-marginal kernel
change; it is not a new general principle that offspring dependence affects
extinction.

\begin{theorem}\label{thm:extinction}
For an exchange-symmetric daughter kernel define
\begin{equation}\label{eq:G}
 G=Q-\diag(b+d)+2\diag(b)L.
\end{equation}
An expected row count vector obeys $m(t)=m(0)e^{tG}$. For two states let
\begin{equation}\label{eq:Kchange}
 K_i^{+}=K_i+\varepsilon_i(1,-1,-1,1),\qquad\varepsilon_i\ge0,
\end{equation}
with all entries nonnegative. All mean trajectories are unchanged, but for every
initial type and $t\ge0$,
\begin{equation}\label{eq:extorder}
                F_i^{+}(t,0)\ge F_i(t,0).
\end{equation}
The same weak ordering holds for eventual extinction. If $d_S\ne d_T$ and
$b_i\varepsilon_i>0$, then
\begin{equation}\label{eq:strict}
 F_i^{+}(t,0)-F_i(t,0)=\frac{b_i\varepsilon_i(d_S-d_T)^2}{3}t^3+O(t^4),
\end{equation}
so the finite-time inequality is strict for sufficiently small positive $t$.
\end{theorem}
\begin{proof}
A type-$i$ event contributes $e_j-e_i$ for a switch, $-e_i$ for death, and
$e_j+e_k-e_i$ for division. Taking conditional expectations gives row $i$ of
\eqref{eq:G}. Yule domination gives finite first moments and justifies the
stopped generator calculation and its limit, so $m'=mG$. The perturbation
\eqref{eq:Kchange} preserves row mass and each daughter marginal, hence $G$.

The offspring PGF difference is the exact polynomial identity
\begin{equation}\label{eq:square}
 \sum_{jk}(K_i^{+}-K_i)(j,k)x_jx_k=\varepsilon_i(x_S-x_T)^2\ge0.
\end{equation}
Let $f,f^+$ be the right sides of \eqref{eq:pgf}. The unit cube is invariant: on
a face $x_i=0$ the field points inward, and on $x_i=1$ the switch terms and
$b_i(\sum K_ix_jx_k-1)$ are nonpositive. Both fields are cooperative, because for
$j\ne i$ their partial derivatives are $q_{ij}$ plus a nonnegative derivative of
the offspring polynomial on the cube. Equation~\eqref{eq:square} implies
$f^+(x)\ge f(x)$ there.

Comparison follows by a strict-supersolution argument. On any fixed finite
horizon, add $\eta e^{Ct}\one$ to the new solution, with $C$ larger than a
Lipschitz constant of $f$ on a neighbouring positive cube. For sufficiently
small $\eta>0$ this is a strict supersolution of the old ODE. At a first contact
with the old solution, the touching coordinate is equal and the other
coordinates are no smaller, contradicting cooperativity and the strict
differential inequality. Letting $\eta$ decrease to zero proves
\eqref{eq:extorder}. Extinction-by-time events increase with $t$; their limits
yield the eventual ordering.

At zero, $F_i'(0,0)=d_i$, so $(F_S^+-F_T^+)^2=(d_S-d_T)^2t^2+O(t^3)$. The
difference equation is the old field difference plus this forcing multiplied by
$b_i\varepsilon_i$. The old field is Lipschitz, so Gronwall's inequality first
gives $F^+-F=O(t^3)$. Its old-field contribution to the derivative is then
$O(t^3)$. Integrating the leading $t^2$ forcing proves \eqref{eq:strict}.
\end{proof}

If $b_i=b$ and $d_i=d$ are state independent, the total count process is a
birth--death process independent of $Q,K,\pi$. Common demographic clocks
therefore give equality of all unmarked count laws. This is a zero-effect
boundary: increased sister concordance need not strictly increase eventual or
finite-time extinction in every source. Eventual strictness is not proved by the
cubic expansion, since both models may eventually go extinct with probability
one.

\begin{figure}[tb]
\centering\includegraphics[width=\linewidth]{figures/extinction.pdf}
\caption{Model predictions, not biological measurements. Each kernel pair has
exactly the same mean generator, so its two mean curves coincide (left).
Numerical solutions of the branching ODE give extinction (middle) and the
absolute extinction difference (right). The baseline perturbation increases
24-hour extinction by approximately $0.0005464$; the explicitly synthetic
large-contrast source gives approximately $0.05412$. The theorem proves the
ordering, not these floating-point effect sizes. Parameters are specified in the
text and retained with the figure data.}
\label{fig:extinction}
\end{figure}

Figure~\ref{fig:extinction} illustrates Theorem~\ref{thm:extinction} on two
sources. Using \eqref{eq:example} with $q_{ST}=.18$ and
$\varepsilon_S=\varepsilon_T=.05$, the 24-hour extinction increase is
approximately $0.0005464$. It is small and is not presented as practically
important. Guided by the forcing in \eqref{eq:square}, a second synthetic source
uses $q_{ST}=q_{TS}=.01$, $b=(.6,.4)$, $d=(.5,.005)$, $\pi_T=.35$,
independent-daughter base marginals $L_S(T)=.25$, $L_T(T)=.75$, and
$\varepsilon_i=.18$. All kernel entries remain positive. Its 24-hour increase is
approximately $.05412$; mean growth is still identical within the pair. No
compatible biological calibration is known here to justify that larger death
contrast. A long-horizon prediction additionally assumes that the short-assay
source continues to apply.

The implication is specific: expected burden alone can miss an
extinction-relevant component of inheritance. A clinical recommendation would
require further assumptions about intervention effects, toxicity, competition
and the persistence of the fitted source, all outside this paper.

\section{Experimental meaning, model challenges and practical scope}
\subsection{A single biological context, without invented calibration}
Consider the AU565 lineage-imaging architecture studied by Mohammadi et
al.~\citep{mohammadi2022}: drug-treated and control cells are tracked for
cell-cycle phase, division and fate. It is a concrete architectural starting
point, not a calibrated instance of our phenotype model.
Table~\ref{tab:provenance} keeps source provenance separate from synthetic
choices. We do not combine rates from unrelated cell lines to construct a
supposedly calibrated organism.

\begin{table}[tb]
\caption{Parameter provenance for an AU565-inspired assay. Numerical model
inputs are synthetic unless explicitly described as reported.}
\label{tab:provenance}\centering\small
\begin{tabular}{>{\raggedright\arraybackslash}p{.20\linewidth}>{\raggedright\arraybackslash}p{.35\linewidth}>{\raggedright\arraybackslash}p{.35\linewidth}}
\toprule
Quantity & Value or requirement & Status and source\\
\midrule
Context & AU565; lapatinib or gemcitabine lineage imaging & Reported architecture, Mohammadi et al.~\citep{mohammadi2022}, methods\\
Cadence/horizon & Images every 30 minutes; AU565 follow-up 96 hours & Reported; not the instantaneous event detection assumed here\\
Recorded reporter & Nuclear/cell-cycle reporter and terminal fates & Reported; not a calibrated tolerance-state channel\\
States $S,T$ & Predictive phenotypic response classes & Assumed model; no validated AU565 mapping supplied\\
Marker $E$ & $e_S=.1$, $e_T=.85$ & Synthetic; labelled or known-mixture calibration required\\
Rates and $K$ & Equation~\eqref{eq:example}; hours and inverse hours & Synthetic, not estimated from AU565 records\\
Clock/delay & $\lambda\in\{.25,.5,1,2\}$ per hour; known $\tau$ & Design choices; detect interruptions explicitly\\
Independent units & Independent founders/blocks & Required; related cells in one lineage are not iid roots\\
\bottomrule
\end{tabular}
\end{table}

With the marker and rate assumptions met, the experiment estimates both
within-cell switching and division-associated daughter production. For the
synthetic source \eqref{eq:example}, $L_S(T)=.26$ and $L_T(T)=.8$, so the
inheritance contrast is $.54$ and the sister covariances are $.0724,.06$.
Tolerant daughter-production intensities are $\zeta_S=.1144$ and $\zeta_T=.256$
per hour. They are equal across the two example arms by construction; induction
differs by $.8$ per hour. In an experiment neither equality need be assumed.
Selection, inheritance and induction may all coexist, and none of them alone
identifies genetic versus epigenetic causation.

Multiplying every biological rate and $\lambda$ by $c>0$, and dividing
$\tau,L_0$ by $c$, leaves the ideal record law unchanged. Thus ratios such as
switching/division, death/division and delay times event rate govern this design.
Scaling time cannot validate those ratios or supply a missing marker. A
population doubling time, a mean cell-cycle duration and an exponential
single-cell division hazard are not interchangeable parameters.

\subsection{Age dependence and what a phase channel must supply}
Constant hazards are a substantive assumption. If the division time $D$ has
Laplace transform $\varphi(\lambda)=\E e^{-\lambda D}$ and there is no death or
switching, then $\Prob(T<D)=1-\varphi(\lambda)$ and the one-state exponential
source inverse produces an apparent birth rate
\begin{equation}\label{eq:age}
 b_{\rm app}(\lambda)=\lambda\bigl(A_\lambda^{-1}-1\bigr)
  =\lambda\frac{\varphi(\lambda)}{1-\varphi(\lambda)}.
\end{equation}
For $D\sim\mathrm{Erlang}(2,\nu)$, $\varphi(\lambda)=(\nu/(\nu+\lambda))^2$ and
$b_{\rm app}(\lambda)=\nu^2/(2\nu+\lambda)$. At $\nu=1$, the deadline rates
$\lambda=1,2$ give $1/3$ and $1/4$, so no common homogeneous one-state generator
explains both deadline groups. This is a model obstruction, not sampling noise,
and it is experimentally testable: it converts the cross-deadline check
\eqref{eq:identity} into a falsification instrument.

Multiple deadline rates therefore measure aspects of the waiting-time
distribution even when the Markov inverse fails. A continuum of Laplace
transform values determines the waiting-time law under standard uniqueness
conditions; finitely many noisy values do not identify an unrestricted
distribution. Estimating a constrained waiting-time law is a well-posed future
target, and it is not solved here.

A finite-phase extension is available but is not automatic. Expand the state to
(phenotype, cell-cycle phase), specify the phase transitions and the division
reset, and apply Theorem~\ref{thm:resolvent} \emph{only if its full observation
and preparation assumptions hold for the expanded chain}. Three conditions are
easy to lose. First, the calibrated channel must distinguish all expanded
states: if a two-state tolerance channel and a two-phase channel have
independent errors conditional on the expanded state and each has rank two, their
tensor channel has rank four; otherwise the joint channel must be calibrated
directly. A phenotype-only marker has identical columns for the hidden phases
and is not full rank, and FUCCI colour classes need not separate every hidden
Erlang stage. Second, founder support must be positive on every expanded state:
a cohort of newborn cells has zero mass in later phases, violating the positive
founder masses used for row normalization, so an asynchronous calibrated
preparation or a new reachability argument is required. Third, nondividing
phases have $b_i=0$, where only the zero exit mass is recovered. The initial age
distribution matters in its own right: an asynchronously sampled cell has a
residual lifetime distribution that need not equal its newborn lifetime
distribution, and age stratification alone does not make an age-dependent hazard
constant. Internal-state and age-dependent branching models are developed in
\citep{gunnarsson2023}; a phenotype-plus-phase extension should be positioned
against that literature.

Likewise, $F_\tau$ describes delayed \emph{biological state evolution}, not
reporter maturation. If two hidden states have identical emissions and identical
aggregated transition, demographic and daughter laws under every admitted
intervention, they decorate the same observed process: all record laws agree,
while the founder split and hidden inheritance can vary.

\subsection{Destructive readouts and well-level dependence}
Destructive initial staining cannot be applied to the same founder that is then
followed. The initial marker requirement therefore needs a compatible live
readout, a justified alternative preparation design, or a different theorem.
Randomizing fixation times at the well level does not by itself supply the
missing initial measurement or an immediate daughter-state measurement.

If $k$ cells share one random well deadline $T_w$, their outcomes may be
independent conditional on $T_w$ under an idealized source, but they are
generally dependent after marginalizing it. For a category probability $p(T_w)$,
independent wells and conditionally iid cells give
\begin{equation}\label{eq:well}
 \operatorname{Var}(\bar X)
 =\frac{\operatorname{Var}\bigl(p(T_w)\bigr)}{W}
 +\frac{\E\bigl[p(T_w)(1-p(T_w))\bigr]}{Wk}.
\end{equation}
Increasing cells per well does not remove the first term. One safe concentration
construction treats each well's bounded average as a single independent
observation and applies the bounded-difference argument across the $W$ wells; a
richer conditional-on-time likelihood can use within-well information more
efficiently, but it is not the iid-founder multinomial analysed above. Arbitrary
shared well effects can create clone-like correlations, and treatment
randomization alone does not turn a mixture over environments into a homogeneous
Markov source. Use measured or controlled blocks and replicate independent units
at the level on which errors are computed.

Throughput is a real constraint, but no universal hardware ceiling follows from a
platform name. What must be quantified for a specific instrument is the number of
independent preparations, imaging fields, tracking failure rate, calibration
acquisition, total cell-hours and the actual observed state channels.

\begin{figure}[tb]
\centering\includegraphics[width=\linewidth]{figures/delay.pdf}
\caption{Delay conditioning for the treated source in \eqref{eq:example}.
Singular values are numerical evaluations of the exact channel \eqref{eq:delay};
the pair-channel value is their square. The dashed $0.01$ spectral-error level
illustrates where a norm perturbation bound alone cannot certify pair-channel
rank; it is not an empirical calibration error and not a theorem that
information vanishes at that delay. Right: additional daughter cell-hours per
founder are bounded by $2\tau y_\lambda$, with $y_1\simeq.14624$. Founder
observation averages $.79690$ hours in this arm. Long delay retains structural
rank but can destroy practical precision.}
\label{fig:delay}
\end{figure}

\subsection{A usable design decision}
Table~\ref{tab:design} distinguishes identification from a ready assay. A cap of
$14/\lambda$ hours has reserve below $10^{-6}$; known finite caps can instead be
modelled exactly. Daughter monitoring may continue for $\tau$ after founder
division, so elapsed follow-up is at most $L_0+\tau$ under a cap. The
independently labelled calibration samples in the worked example total 200,000;
their acquisition cost is additional to founder observation time.

\begin{table}[tb]
\caption{Design requirements and failure signs. Sample counts are independent
roots or labelled calibration readouts, not descendant counts.}
\label{tab:design}\centering\small
\begin{tabular}{>{\raggedright\arraybackslash}p{.22\linewidth}>{\raggedright\arraybackslash}p{.43\linewidth}>{\raggedright\arraybackslash}p{.25\linewidth}}
\toprule
Design or resource & What it supplies and costs & Failure sign / limit\\
\midrule
One binary endpoint & One count-threshold outcome & Preparation/rate alias, Proposition~\ref{prop:alias}\\
Randomized stopped record & Initial marker, event class, endpoint or joint daughters; mean founder time $\le1/\lambda$ & Missing denominators or marker rank\\
Six-category record & Switching target only, no daughter phenotypes, Proposition~\ref{prop:six} & Cannot recover $b,d,K$ or predict extinction\\
Monitored delay & Each daughter: marker at $\tau$ or interruption; at most $2\tau y_\lambda$ extra cell-hours/root & Missed interruptions, unknown reporter memory\\
Two arms and calibration & Worked example: 50,000 roots/arm, 100,000 labelled reads/state & Report unresolved if the determinant bound reaches zero\\
Prospective decision & Narrow class of \S4.6: $1.5\times10^5$ roots/arm and reads/state & Sufficient for that class only; not a lower bound\\
Calibration material & Pure states, or two known mixtures, Proposition~\ref{prop:mixcal} & Circular if gated on the same marker\\
Division information & At $\lambda=1$, treated division yield $.14624$; needed for $K$ & Rare divisions weaken kernel recovery\\
Physical cap & $L_0=14/\lambda$, reserve $<10^{-6}$ in each feature & Treating a capped kernel as an exact resolvent\\
Multiple deadline rates & Intersect common-source confidence sets & Incompatible $H,B,\pi$; cause not uniquely identified\\
Extinction prediction & Identified joint $K$ plus continued source validity & Extrapolation beyond validated times/conditions\\
\bottomrule
\end{tabular}
\end{table}

For an explicit reporter failure, suppose daughter interruption hazards are
state independent and the delayed reporter has identical columns for the two
birth states. The whole daughter channel then has identical columns. Any two
kernels with the same row mass, including different sister covariances, give the
same delayed pair law. Founder data can still recover $H,b,d$, but cannot repair
that rank loss in the daughter readout. Thus an uncalibrated maturation channel
cannot be substituted for the full-rank channel in Theorem~\ref{thm:delay}.

For every limitation there is a specific observation that would address it: an
initial or live state channel for preparation; labelled or independently
characterized mixtures for calibration; demographic flags for denominator
integrity; joint daughters for $K$; independent blocks for well effects; phase
information or waiting-time profiles for age dependence. More uninformative
replicates do not resolve an exact alias.

The experimental conclusion is conditional but specific: with calibrated state
channels, intact denominators and the stated stopping rule, switching and
daughter production are separately recoverable, and joint daughter observations
support an extinction consequence invisible to mean growth. Without those
channels, the paper identifies concrete ambiguities rather than promising
retrospective recovery from data that never measured them.

\appendix
\section{Numerical certification and implementation details}\label{app:numerics}
The exact rational population law for a randomized deadline is particularly
convenient. Rational $\lambda,H,B,E,\pi$ give rational matrices
$A_\lambda,D_\lambda,J,V$; no exponential approximation is used for those laws.
Directed checks include a three-state cycle with
\[
 H=\begin{bmatrix}-3&2&0\\0&-3&2\\2&0&-3\end{bmatrix},\quad
 \lambda=\tfrac32,\quad
 A_\lambda=\frac1{665}\begin{bmatrix}243&108&48\\48&243&108\\108&48&243\end{bmatrix}.
\]
The spectrum of $H$ is $-1,-4\pm i\sqrt3$. Rational matrix multiplication
verifies \eqref{eq:inverse} exactly. Reducible and stiff-rate checks also
include a zero eigenvalue, showing why inversion of $H$ itself is unnecessary.
These examples validate implementation conventions; the integral proof of
Theorem~\ref{thm:resolvent} supplies the general result.

Interval operations are endpoint inclusion operations on exact fractions. For a
positive rational $x$, the upward square-root approximation is
$(\lfloor10^{12}\sqrt{x}\rfloor+1)/10^{12}$, obtained by integer square root of
the scaled numerator and denominator. Matrix products are finite sums of
interval products. A nonpositive lower determinant or marker contrast leads to
an unbounded result, never an arbitrarily chosen denominator floor. For the
prospective boxes, interval inversion of the two-by-two $\lambda I-H$ precedes
the observed-law and target maps. This enclosure need not be sharp; its
inclusion property and the explicit threshold separation are the certificate. No
finite grid is substituted for a continuous parameter box.

Chernoff sets require a certified enclosure of $\{q:n\kl(\hat p\Vert q)\le L\}$.
We locate each endpoint by floating-point bisection, move it outward to a
rational with denominator $10^{9}$, and then \emph{verify} that the candidate
lies outside the set by evaluating a rational lower bound on
$\kl(\hat p\Vert\cdot)$. That lower bound uses two-sided rational bounds on the
logarithm: $\log x=2\operatorname{artanh}\frac{x-1}{x+1}$ with the odd series
truncated after a fixed number of terms and the remainder bounded by the
geometric tail
$|z|^{2N+1}\bigl((2N+1)(1-z^2)\bigr)^{-1}$, preceded by exact halving of the
argument into $[2/3,4/3]$ and a rational enclosure of $\log2$ obtained from the
same series at $z=1/3$. Since $\kl$ is a positive combination of two logarithms
with known signs, term-wise lower bounds give a sound lower bound on $\kl$. A
separate audit re-evaluates every certified endpoint at 200 decimal digits and
confirms that it lies outside the exact Chernoff set.

The script \texttt{check\_paper.py} distributed with the source replays every
constant printed above: the retained dataset regenerated from its seed, the
six-category totals, the soundness of the logarithm and $\kl$ bounds, the
containment of each Chernoff set in its Hoeffding radius, the four rows of
Table~\ref{tab:routes} with their width reductions, the unresolved small-sample
example, the exact resolvent checks, the prospective budgets, the delay channel
values, the Erlang obstruction, the extinction effect sizes and the Monte Carlo
diagnostic. All 86 checks pass.

\section{Reproducibility and proof status}
The chronological construction, the observation laws, the resolvent integral,
matrix-log uniqueness, confidence coverage and the cooperative comparison are
conventional mathematical proofs in this paper. Lean verifies a declared finite
algebra subset: resolvent cancellation, the rational two-state target identity,
rectangular delayed-channel cancellation, the offspring polynomial identity, and
--- new in this version --- the cancelled-determinant identity \eqref{eq:S1} in
the form of its two marker-block lemmas, the known-mixture calibration formulas
\eqref{eq:S2}, the contamination bound of Remark~\ref{rem:contam}, and the
six-category subtraction identity of Proposition~\ref{prop:six}. Successful
compilation of that subset does not constitute a compiled proof of the
stochastic theorems, and this paper must not be described as machine verified.
In particular, Lemma~\ref{lem:chernoff}, Theorem~\ref{thm:CI} and
Proposition~\ref{prop:mixture} are conventional proofs.

The pinned environment uses Lean 4.30.0 and its frozen Mathlib dependency. All
five modules, including the new
\texttt{PhenotypeIdentification/CancelledTarget.lean}, compile with warnings
promoted to errors, exit code zero and empty output. An axiom probe confirms
that each of the six new declarations depends only on \texttt{propext},
\texttt{Classical.choice} and \texttt{Quot.sound}; no custom mathematical axiom
and no incomplete proof is admitted. The script \texttt{lean\_verify.py}
reproduces the compilation and the probe and writes
\texttt{lean\_receipt.json} with exit codes and source hashes.

The source package is self-contained: \texttt{model.py} holds the synthetic
source, the stopped observation law and the branching ODE; \texttt{certkit.py}
holds exact interval arithmetic, the rational logarithm bounds, the Chernoff
sets and the two target evaluators; \texttt{prospective\_boxes.py} holds the
class-level decision; \texttt{make\_figures.py} regenerates the four vector
figures; and \texttt{check\_paper.py} replays every number. Source-record
simulation uses seed 20260922; the 300-repetition coverage and power diagnostic
regenerates the calibration sample and both arms each time. Exact binomial
intervals describe Monte Carlo uncertainty only. Extinction plots use numerical
ODE solutions, not interval-certified effect magnitudes.

All examples are synthetic. No experimental data were fitted, no empirical rate
ranges are claimed, and no clinical conclusion is inferred. Author information,
affiliations and a data and code availability statement are to be completed
before submission.

\bibliographystyle{plainnat}
\bibliography{refs}

\end{document}
